% 整数（综述）
% license CCBYSA3
% type Wiki

本文根据 CC-BY-SA 协议转载翻译自维基百科\href{https://en.wikipedia.org/wiki/Integer}{相关文章}

\begin{figure}[ht]
\centering
\includegraphics[width=10cm]{./figures/f21fcdea866249ee.png}
\caption{排列在数轴上的整数} \label{fig_ZS1_1}
\end{figure}
整数是指 零（0）、正自然数（1, 2, 3, …），或 正自然数的相反数（−1, −2, −3, …）。\(^\text{[1]}\)正自然数的相反数（或加法逆元）称为负整数。\(^\text{[2]}\)所有整数构成的集合通常记作 粗体 \(\mathbf{Z}\) 或 黑板粗体：$\mathbb{Z}$.\(^\text{[3][4]}\)  

自然数集合  $\mathbb{N}$是整数集合$\mathbb{Z}$的子集，而 $\mathbb{Z}$ 又是有理数集合$\mathbb{Q}$的子集，而 $\mathbb{Q}$ 本身是实数集合$\mathbb{R}$的子集。\(^\text{[a]}\)与自然数集合类似，整数集合$\mathbb{Z}$是可数无限集。整数可以被视为一种实数，其表示中不含分数部分。例如：$21, 4, 0, -2048$ 都是整数；而 $9.75, 5+\tfrac{1}{2}, \tfrac{5}{4}, \sqrt{2}$ 都不是整数。\(^\text{[7]}\)

整数构成了包含自然数的最小的群与最小的环。在代数数论中，整数有时被称为有理整数，以区别于更一般的代数整数。实际上，（有理）整数就是那些既是代数整数、又是有理数的数。
\subsection{历史}
integer一词源自拉丁语 integer，意为“完整的”或（字面意义上）“未触碰的”，由 in（“不”）+ tangere（“触碰”）组成。英语单词 entire 也来自同一词源，经由法语 entier（既表示“完整的”，也表示“整数”）。[8]在历史上，该词曾被用于表示1 的倍数[9][10]，或者表示带分数的整数部分。[11][12] 当时仅考虑正整数，因此该词与自然数同义。随着负数的实用性被逐渐认识到，整数的定义才扩展到包括负数。[13]例如，莱昂哈德·欧拉在其 1765 年的《代数学原理》中，就将整数定义为既包含正数也包含负数。[14]

“整数集合”这一说法在 19 世纪末之前并未使用。随着格奥尔格·康托尔引入无限集合与集合论，该说法才出现。用字母 **Z** 表示整数集合，来源于德语单词 *Zahlen*（“数”）[3][4]，这一记号常被归功于大卫·希尔伯特。[15]该符号在教材中最早出现于 1947 年由布尔巴基学派（Nicolas Bourbaki 集体笔名）编写的《代数学》中。[3][16] 但该记号并未立即被普遍接受。例如，有的教材使用字母J，[17] 而 1960 年的一篇论文用 **Z** 表示非负整数。[18] 然而，到 1961 年，现代代数学教材中普遍使用 Z 表示包含正负整数的集合。[19]

符号 $\mathbb{Z}$ 常带上标或下标以表示不同的整数集合，不同作者用法可能不同：

* $\mathbb{Z}^{+}, \mathbb{Z}_{+}, \mathbb{Z}^{>}$：正整数
* $\mathbb{Z}^{0+}, \mathbb{Z}^{\geq}$：非负整数
* $\mathbb{Z}^{\neq}$：非零整数
* $\mathbb{Z}^{*}$：有些作者表示非零整数，有些表示非负整数，还有些表示 $\{-1, 1\}$（即 $\mathbb{Z}$ 的单位元群）
* $\mathbb{Z}_p$：可能表示 **模 $p$ 的整数集合**（即整数的同余类集合），也可能表示 **$p$-进制整数集合**。\[20]\[21]

直到 1950 年代初期，**whole numbers（全体数）** 一词还被认为与整数同义。\[22]\[23]\[24]
在 1950 年代末期，作为“新数学运动（New Math movement）”的一部分，美国小学教师开始教授 **whole numbers** 指自然数（不含负数），而 **integer** 则包括负数。\[26]\[27]
直到今天，**whole numbers** 一词仍然存在歧义。\[28]

---

要不要我帮你整理一个 **整数符号用法对照表（$\mathbb{Z}^+ , \mathbb{Z}^* , \mathbb{Z}_p$ 等）**，让这些记号更直观清晰？
